\section{Summary of Contributions}
While the shift from centralized to DFL successfully eliminates the risks of a single point of failure and centralized data monopolies, it places the burden of security squarely on the network's structural topology. This thesis systematically evaluated the structural vulnerabilities of these peer-to-peer environments, utilizing a continuous backdoor payload strictly as an experimental mechanism to observe how network geometry dictates parameter flow.

At its core, this research set out to answer a critical question: how does the physical layout of a network influence the survival and spread of a localized backdoor payload? By deploying a targeted threat across diverse architectural constraints, ranging from dense Fully Connected clusters to realistic, community-based Stochastic Block Models, this thesis provides empirical proof that network geometry acts as the fundamental variable governing adversarial propagation. 

To enable a fair and mathematically sound evaluation, this research formulated and validated a topological boosting heuristic, $\alpha=d+1$. By scaling the malicious parameter updates proportionally to the attacker's local topological degree, the adversary was able to perfectly neutralize the natural dilution caused by neighborhood aggregation. Rather than focusing on maximizing raw attack damage, this heuristic was utilized to maintain strict metric stability. The attacker consistently maintained a Clean Accuracy indistinguishable from the honest majority, successfully preserving the experimental baseline while exposing the underlying structural weaknesses of the graph. 

The strategic placement of the adversary across these diverse structures yielded several critical insights into network resilience:

\begin{itemize}
    \item \textbf{Density Accelerates Infection:} In a maximally dense Fully Connected topology, the absence of communication bottlenecks allows a continuous injection to easily overcome local dilution, driving a rapid and system-wide compromise.
    \item \textbf{Sparsity Delays Propagation:} Constrained topologies, such as the Ring, naturally throttle the infection. The physical necessity of sequential, multi-hop communication slows the payload's progression, demonstrating that structural distance natively buys the network crucial time.
    \item \textbf{The Hub Vulnerability:} In a Star topology, centralized aggregation acts as an impenetrable shield against anomalous updates originating from the periphery. However, if the central Hub itself is compromised, the topology becomes exceptionally fragile, granting the attacker an unfiltered broadcast channel to instantly infect the entire network.
    \item \textbf{The Practical Optimum (SBM):} Evaluations using the Stochastic Block Model proved that community boundaries act as natural defenses, making it the ideal architecture for real-world deployment. It provides the rapid local convergence of a dense network, while its sparse gateway links act as natural topological firewalls. Even when the attacker controls the gateway node itself, the target community functions as a consensus sink, structurally washing out the malicious parameters and preventing global corruption.
\end{itemize}

Finally, by replicating these dynamics on the NEU-DET industrial dataset, this research confirmed that these structural vulnerabilities are inherent to the decentralized architecture, completely independent of the underlying complexity of the training data. Crucially, the industrial validation revealed that in highly complex environments, community boundaries shift from absolute firewalls to powerful temporal delays.

By trapping the payload for hundreds of rounds, the topology naturally buys the network the critical time needed to run diagnostics or deploy active intervention protocols.

\section{Limitations of the Study}
Like any controlled experiment, this research required abstracting certain real-world complexities to precisely isolate the variables under study. The primary limitations of this methodology are:

\begin{itemize}
    \item \textbf{Scale of the Network:} To precisely observe the step-by-step mathematical dilution of the malicious payload, the experiments were conducted on an 8-node network. While this scale perfectly demonstrates the mechanics of peer-to-peer architectures, enterprise-level edge computing networks may consist of thousands of participating devices, which could alter the speed of global consensus.
    \item \textbf{Static Graph Architecture:} The evaluated network topologies were strictly static. In real-world environments, peer-to-peer networks are highly dynamic and characterized by churn, where nodes frequently disconnect, rejoin, or establish new lateral connections mid-training due to bandwidth or battery constraints.
    \item \textbf{Uniform Aggregation Assumption:} The formulation of the topological boosting heuristic fundamentally relies on the assumption of uniform neighborhood averaging, where honest nodes weight all incoming parameters equally. If a network were to employ dynamic weighting or trust-based aggregation algorithms, the multiplier would no longer perfectly cancel out the receiver's dilution. In such environments, an attacker would have to dynamically estimate and adapt to the specific weights used by their neighbors.
    \item \textbf{IID Data Distribution:} The simulations utilized an Independent and Identically Distributed (IID) partitioning of both the MNIST and NEU-DET datasets. While this successfully controlled for data-driven variance, practical decentralized networks often suffer from highly heterogeneous (Non-IID) data distributions across different clients.
\end{itemize}

\section{Future Lines of Research}
Building directly upon the findings and limitations of this thesis, several exciting avenues for future research emerge in the domain of decentralized AI security:

\begin{itemize}
    \item \textbf{Evaluating Non-IID Environments:} Future work must investigate how the $\alpha=d+1$ topological boosting heuristic performs under severe data heterogeneity. It remains to be seen whether Non-IID data distributions act as an additional layer of natural dilution against backdoor payloads, or if the resulting chaotic loss landscape actually makes it easier for anomalous parameters to propagate.
    \item \textbf{Dynamic and Evolving Topologies:} Simulating the continuous backdoor attack on dynamic graphs is a crucial next step. Specifically, researchers should explore how community-based topological firewalls behave under network churn. For example, if a gateway node disconnects and the network dynamically reroutes its connections, it could inadvertently bridge an isolated attacker directly to the broader network.
    \item \textbf{Topology-Aware Defense Mechanisms:} Finally, while this thesis proves that structure alone can significantly delay continuous injections, future systems must develop active, topology-aware defenses. The design of Byzantine-robust aggregation algorithms tailored specifically for DFL, where nodes dynamically adjust their aggregation weights based on the trust or historical variance of their structural neighbors, remains a highly promising frontier for securing peer-to-peer AI infrastructure.
\end{itemize}